\documentclass[11pt]{article}
\usepackage{amsmath,amssymb,amsthm}

\title{Maximal Integral Curves and Flow Extensions}
\author{UlamLib seed note}
\date{}

\newtheorem{definition}{Definition}[section]
\newtheorem{lemma}[definition]{Lemma}
\newtheorem{theorem}[definition]{Theorem}
\newtheorem{corollary}[definition]{Corollary}

\begin{document}
\maketitle

\section{Scope}

This note targets a clean global ODE layer on top of local integral-curve
machinery. The intended Lean destination is
\texttt{UlamLib/ODE/MaximalIntegralCurves.lean}.

\paragraph{Target Lean declarations.}
This note should eventually feed declarations with names close to
\texttt{IsRightExtendable}, \texttt{IntegralCurveExtension},
\texttt{MaximalIntegralCurve}, and \texttt{maximalCurve\_unique}.

\paragraph{Target Lean file family.}
The track is intentionally split across
\texttt{UlamLib/ODE/Extendability.lean},
\texttt{UlamLib/ODE/MaximalIntegralCurve.lean}, and
\texttt{UlamLib/ODE/LinearSystems.lean}.

\section{Ambient assumptions}

Fix a smooth manifold $M$ without boundary, or restrict attention to initial
points in the interior when boundary issues would matter. Let $V$ be a smooth
vector field on $M$, and let $I,J \subseteq \mathbb{R}$ be open intervals.
Every curve in this note is assumed to be a smooth integral curve of $V$ on
its stated domain.

\section{Declaration blueprint}

\begin{enumerate}
\item \texttt{IntegralCurveExtension}: data of one curve extending another.
\item \texttt{IsRightExtendable}: extendability predicate at a right endpoint.
\item \texttt{IsLeftExtendable}: left-endpoint variant.
\item \texttt{MaximalIntegralCurve}: a curve together with its maximal open
interval of existence.
\item \texttt{compatibleOnOverlap}: compatibility relation used for gluing.
\end{enumerate}

\section{Core definitions}

\begin{definition}[Integral curve]\label{def:integral-curve}
Let $M$ be a smooth manifold and let $V$ be a smooth vector field on $M$. An
integral curve of $V$ is a smooth map $\gamma : I \to M$ on an interval
$I \subseteq \mathbb{R}$ such that $\gamma'(t) = V(\gamma(t))$ for all
$t \in I$.
\end{definition}

\begin{definition}[Right extendability]\label{def:right-extendability}
A local integral curve is extendable to the right endpoint of its interval if
there exists a strictly larger interval and an integral curve on that larger
interval that restricts to the original curve.
\end{definition}

\begin{definition}[Maximal integral curve]\label{def:maximal-integral-curve}
A maximal integral curve is an integral curve together with an open interval of
existence that admits no further extension.
\end{definition}

\section{Seed theorem cluster}

\begin{lemma}[Overlap uniqueness]\label{lem:overlap-uniqueness}
If two local integral curves of the same smooth vector field agree at one point
in the intersection of their domains, then they agree on every connected
component of the overlap containing that point.
\end{lemma}

\begin{lemma}[Patching on overlap]\label{lem:patching-on-overlap}
Compatible local integral curves can be glued along overlapping intervals to
produce a larger local integral curve.
\end{lemma}

\begin{theorem}[Uniqueness of maximal curves]\label{thm:uniqueness-of-maximal-curves}
Two maximal integral curves with the same initial condition agree on the
intersection of their domains.
\end{theorem}

\begin{corollary}[Existence from extendability data]
\label{cor:existence-from-extendability}
If every local integral curve through a point can be extended past each finite
endpoint, then there exists a maximal integral curve through that point.
\end{corollary}

\section{Formalization order}

\begin{enumerate}
\item Package extension data and endpoint extendability predicates.
\item Prove overlap uniqueness and the gluing lemma.
\item Define maximal integral curves as bundled objects.
\item Prove maximal-curve uniqueness from overlap uniqueness.
\item Add continuation criteria for bounded vector fields or compact manifolds.
\item Connect constant and linear systems to the matrix exponential API.
\end{enumerate}

\end{document}
